\section{哈密顿同胚群}
\label{sec:hameomorphism-group}
\stfhIndex{I159}{ha^mi^dun^tong^pei^qun}{哈密顿同胚群（Hamiltonian homeomorphism group）}

(OhM07) 的作者们引入了如下给出的\emph{哈密顿同胚}概念：
\begin{equation}
\begin{split}
 \operatorname{Hameo}(M,\omega)
 =\bigl\{h\in\operatorname{Homeo}(M)\mathrel{\big|}\;&
 h=\overline{\operatorname{ev}}_1(\lambda),\\
 &\lambda\in\mathcal P_\infty^{\mathrm{ham}}
 (\operatorname{Sympeo}(M,\omega),\operatorname{id})\bigr\}.
\end{split}
 \label{eq:hameomorphism-group}
\end{equation}
这个概念是在哈密顿路径的 $L^{(1,\infty)}$-拓扑背景下给出的。下面的正规性是在 (OhM07) 中于
$L^{(1,\infty)}$ 背景下证明的。

\begin{theorem}\label{thm:hameomorphism-normal-path-connected}
$\operatorname{Hameo}(M,\omega)$ 是
$\operatorname{Sympeo}(M,\omega)$ 的道路连通正规子群（关于
$\operatorname{Homeo}(M)$ 的子空间拓扑）。
\end{theorem}

(OhM07) 中给出的这个定理之证明相当精细，主要原因是哈密顿量的
$L^{(1,\infty)}$-极限与哈密顿路径的 $C^0$-极限之间的相互作用涉及一些问题。另一方面，可以重复前一节给出的构造，并在
$[0,1]\times M$ 上的 $C^0$ 背景中进行；这比在
$L^{(1,\infty)}$ 背景中更为简单。这样便得到
$\operatorname{Sympeo}(M,\omega)$ 的两个先验上不同的正规子群，暂记为
$\operatorname{Hameo}_{(1,\infty)}(M,\omega)$ 与
$\operatorname{Hameo}_\infty(M,\omega)$。

但是，穆勒 (Mue08) 证明了下面的定理。

\begin{theorem}[穆勒]\label{thm:hameomorphism-topologies-agree}
有
\[
 \operatorname{Hameo}_{(1,\infty)}(M,\omega)
 =\operatorname{Hameo}_\infty(M,\omega).
\]
\end{theorem}

在此基础上，从现在起，可以把任一个哈密顿同胚群都简写成
$\operatorname{Hameo}(M,\omega)$。由于上述定理 6.4.2，在定理
6.4.1 的证明中，可以使用作为光滑哈密顿量之极限的
$C^0$ 对象。于是，正规性与道路连通性的证明变得非常简单。

\begin{proof}[定理 6.4.1 的证明]
本证明采用定理 6.4.2 给出的、以
$\operatorname{Hameo}_\infty(M,\omega)$ 表示 $\operatorname{Hameo}(M,\omega)$ 的描述。设
\[
 h\in\operatorname{Hameo}_\infty(M,\omega),
 \qquad
 g\in\operatorname{Sympeo}(M,\omega).
\]
由这两个群的定义，存在光滑哈密顿量序列
$H_i$ 以及序列 $\psi_i\in\Symp(M,\omega)$，使得
$(\phi_{H_i},H_i)$ 与 $\psi_i$ 在 $C^0$-拓扑中收敛，并且
\[
 \lim_{i\to\infty}\phi_{H_i}^1=h,
 \qquad
 \lim_{i\to\infty}\psi_i=g.
\]
显然，$H_i\circ\psi_i\to H\circ g$，而且
\[
 \psi_i^{-1}\phi_{H_i}\psi_i
 =\phi_{H_i\circ\psi_i}
\]
在 $C^0$-拓扑中收敛；这首先蕴含 $H\circ g$ 是拓扑哈密顿量，继而蕴含
$H_i\circ\psi_i$ 是它的一个逼近序列。此外，
\[
 \psi_i^{-1}\phi_{H_i}^1\psi_i
 =\operatorname{ev}_1(\psi_i^{-1}\phi_{H_i}\psi_i)
\]
且
\[
 \lim_{i\to\infty}\psi_i^{-1}\phi_{H_i}^1\psi_i
 =g^{-1}hg.
\]
所以依定义，已经证明
\[
 g^{-1}hg\in\operatorname{Hameo}_\infty(M,\omega).
\]
这就证明了 $\operatorname{Hameo}_\infty(M,\omega)$ 在
$\operatorname{Sympeo}(M,\omega)$ 中的正规性。道路连通性也立即可见：上述证明表明，对任意
\[
 h\in\operatorname{Hameo}_\infty(M,\omega),
\]
都存在一条完全位于 $\operatorname{Hameo}_\infty(M,\omega)$ 内、且在
$\operatorname{Sympeo}(M,\omega)$ 中连续的路径，把 $h$ 连接到恒等元。
\end{proof}

与区间 $[0,1]$ 的情形类似，当 $b>a$ 时，可以把
$[a,b]$ 上的拓扑哈密顿路径
\[
 \lambda:[a,b]\longrightarrow\operatorname{Hameo}(M,\omega)
\]
定义为满足
\begin{equation}
 \lambda\circ(\lambda(a))^{-1}
 \in\mathcal P_{[a,b]}^{\mathrm{ham}}
 (\operatorname{Sympeo}(M,\omega),\operatorname{id})
 \label{eq:topological-hamiltonian-path-general-interval}
\end{equation}
的路径，其中
$\mathcal P_{[a,b]}^{\mathrm{ham}}
(\operatorname{Sympeo}(M,\omega),\operatorname{id})$
的定义与
$\mathcal P^{\mathrm{ham}}
(\operatorname{Sympeo}(M,\omega),\operatorname{id})$
相同，只是以 $[a,b]$ 代替 $[0,1]$。

我们相信下面的猜想成立。

\begin{conjecture}\label{conj:hamiltonian-continuous-functions-proper}
子集 $C_{\mathrm{ham}}^0(M,\omega)$ 是 $C^0(M)$ 的真子集。
\end{conjecture}

利用拓扑哈密顿量的唯一性结果 (BS10)，可以把霍费尔长度函数的定义延拓到拓扑哈密顿路径，并通过延拓
$\Ham(M,\omega)$ 上的内蕴距离函数，在
$\operatorname{Hameo}(M,\omega)$ 上定义其相伴内蕴距离函数。更多细节以及其他相关讨论见 (Oh10)。

下面的定理是 $\Ham(M,\omega)$ 上霍费尔范数之著名定理 (H90) 的拓扑哈密顿对应物。

\begin{theorem}\label{thm:extended-hofer-norm-hameomorphisms}
设 $g,h\in\operatorname{Hameo}(M,\omega)$。则延拓的霍费尔范数函数
\[
 \|\mathord\cdot\|:\operatorname{Hameo}(M,\omega)
 \longrightarrow\mathbb R_+
\]
在哈密顿拓扑中连续，并满足下列性质：
\begin{enumerate}[label=(\arabic*)]
 \item （对称性）$\|g\|=\|g^{-1}\|$；
 \item （双不变性）$\|gh\|=\|hg\|$；
 \item （三角不等式）$\|gh\|\leq\|g\|+\|h\|$；
 \item （辛不变性）对任意
 $\psi\in\operatorname{Sympeo}(M,\omega)$，有
 $\|\psi^{-1}g\psi\|=\|g\|$；
 \item （非退化性）$g=\operatorname{id}$ 当且仅当 $\|g\|=0$。
\end{enumerate}
\end{theorem}

然而，穆勒 (Mue08) 提出了如下问题。

\begin{question}[穆勒]\label{ques:hofer-norm-restriction-hameo}
分别以 $\|\mathord\cdot\|_{\mathrm{Ham}}$ 与
$\|\mathord\cdot\|_{\mathrm{Hameo}}$ 表示
$\Ham(M,\omega)$ 上的霍费尔范数以及
$\operatorname{Hameo}(M,\omega)$ 上的延拓霍费尔范数。设
$\phi\in\Ham(M,\omega)\subset\operatorname{Hameo}(M,\omega)$。恒等式
\[
 \|\phi\|_{\mathrm{Ham}}=\|\phi\|_{\mathrm{Hameo}}
\]
是否一般成立？
\end{question}

从第 1 章 (Gr88) 的观点，把这个问题联系于由一般度量空间
$(X,d)$ 出发构造\emph{道路度量空间} $(X,d_\ell)$ 的一般方法，可以为它提供一些启示。

下面的猜想是理解 $C^0$ 层次上哈密顿动力学的重要动机之一。

\begin{conjecture}\label{conj:calabi-homomorphism-hameomorphisms}
卡拉比同态
\[
 \operatorname{Cal}:\Ham(D^2,\partial D^2)\longrightarrow\mathbb R
\]
可延拓为同态
\[
 \overline{\operatorname{Cal}}:
 \operatorname{Hameo}(D^2,\partial D^2)\longrightarrow\mathbb R,
\]
并且该延拓在哈密顿拓扑中连续。
\end{conjecture}

(Oh10) 证明了：这个猜想若成立，再结合磨光定理 (Oh06b)、(Sik07)，将蕴含
$\operatorname{Hameo}(D^2,\partial D^2)$ 是
$\operatorname{Homeo}^\Omega(D^2,\partial D^2)$ 的真子，从而导出
$\operatorname{Homeo}^\Omega(D^2,\partial D^2)$ 非单性的一个证明。在第 4 部中，等到说明谱不变量的构造以及恩托夫--波尔捷罗维奇拟同态 (EnP03) 的概念之后，我们将进一步讨论这个猜想。
